\section{Flujo de pesos, geometría de información y dualidad de Legendre}
\label{sec:convex-refinement}

La deformación de los intervalos admite una descripción variacional
exacta. Sean \(d_j=K_j/Q>0\), \(\sum_jd_j=1\), y \(u_j\) las
coordenadas de la dirección \eqref{eq:u-main}. Para \(s\in\RR\),
se definen
\begin{equation}
 Z(s)=\sum_{j=1}^{12}d_j e^{su_j},\qquad
 \psi(s)=\log Z(s),\qquad
 d_j(s)=\frac{d_j e^{su_j}}{Z(s)}.
 \label{eq:convex-partition}
\end{equation}
Todos los argumentos de la exponencial son adimensionales. El parámetro
\(s\) es la coordenada de esta deformación geométrica. La asignación de
una duración física pertenece a un lector temporal declarado.

\begin{theorem}[Convexidad estricta y coordenada dual]
La función \(\psi\) es analítica y estrictamente convexa sobre \(\RR\).
Si \(u_- =\min_j u_j\) y \(u_+=\max_j u_j\), la aplicación
\(v=\psi'(s)\) es un difeomorfismo de \(\RR\) sobre \((u_-,u_+)\).
Para cada \(v\) en ese intervalo existe una única coordenada \(s(v)\), y
\[
 \psi^*(v)=s(v)v-\psi(s(v)),\qquad
 (\psi^*)'(v)=s(v),\qquad
 (\psi^*)''(v)=\frac1{\psi''(s(v))}>0.
\]
\end{theorem}
\begin{proof}
Las sumas finitas permiten diferenciar término a término. Se obtiene
\begin{align}
 \psi'(s)&=\sum_jd_j(s)u_j=:\overline u(s),\notag\\
 \psi''(s)&=\sum_jd_j(s)(u_j-\overline u(s))^2.
 \label{eq:variance}
\end{align}
Los pesos son estrictamente positivos y los \(u_j\) no son todos iguales;
la varianza es por ello estrictamente positiva. Cuando \(s\) tiende a
\(+\infty\), dividir numerador y denominador de \(\psi'\) por
\(e^{su_+}\) muestra que sólo sobreviven los índices con \(u_j=u_+\).
Así \(\psi'(s)\to u_+\). El argumento con \(u_-\) da el límite
opuesto. La monotonía y el teorema de la función inversa proporcionan la
biyectividad suave. La función \(sv-\psi(s)\) tiene su máximo único
en \(s(v)\); diferenciar el valor máximo demuestra las dos identidades
duales.
\end{proof}

La evolución de cada intervalo es
\begin{equation}
 \dot d_j(s)=d_j(s)\bigl(u_j-\overline u(s)\bigr).
 \label{eq:replicator}
\end{equation}
La suma de las derivadas es cero. Por tanto la longitud total permanece
unitaria mientras se redistribuyen las longitudes internas. Esta conservación
de escala global permite atribuir el cambio de las razones dobles a una
deformación de posiciones y no a una dilatación común.

En el simplex abierto de pesos, la métrica
\(g_d(\xi,\zeta)=\sum_j\xi_j\zeta_j/d_j\), sobre vectores tangentes
de suma cero, es positiva. La función lineal
\(F(d)=\sum_jd_ju_j\) tiene como gradiente para esta métrica
precisamente el campo de \eqref{eq:replicator}, pues
\[
 g_d\bigl(d_j(u_j-\overline u),\xi\bigr)
 =\sum_j(u_j-\overline u)\xi_j=\sum_ju_j\xi_j=dF(\xi).
\]
En consecuencia \(dF/ds=\psi''(s)>0\) a lo largo del flujo. El
carácter ascendente corresponde a esta función y a esta métrica concretas.

\begin{theorem}[Invariancia variacional por subdivisión heredada]
Para cada intervalo \(j\), sean \(r_{j,c}>0\),
\(\sum_c r_{j,c}=1\), los pesos relativos de una subdivisión fija.
Asigne a los sucesores
\[
 d_{j,c}=r_{j,c}d_j,\qquad u_{j,c}=u_j.
\]
La función de partición, su logaritmo, todas sus derivadas y la transformada
de Legendre coinciden exactamente con los de la partición madre.
Además, el flujo conmuta con la suma de sucesores:
\[
 \sum_c d_{j,c}(s)=d_j(s),\qquad
 d_{j,c}(s)=r_{j,c}d_j(s).
\]
\end{theorem}
\begin{proof}
Por agrupación de la suma finita,
\[
 Z_{\rm ref}(s)=\sum_{j,c}r_{j,c}d_j e^{su_j}
 =\sum_j d_j e^{su_j}=Z(s).
\]
Todas las conclusiones se deducen de esta identidad y de dividir cada
peso por el mismo \(Z(s)\). El mismo argumento vale después de cualquier
número finito de subdivisiones. Para una subdivisión numerable, la suma de
términos positivos converge al mismo valor mediante su agrupación por predecesores.
\end{proof}

La hipótesis de dirección heredada tiene contenido verificable. Si un
intervalo con dirección \(a\) se divide en dos mitades de direcciones
\(a+b\) y \(a-b\), su contribución cambia de
\(d e^{sa}\) a \(d e^{sa}\cosh(sb)\). Para \(b\ne0\) la igualdad
de funciones de partición falla aunque el peso total en \(s=0\) sea
el mismo. Así, la conservación geométrica y variacional exige conservar
también la información direccional pertinente. La identidad demostrada
se aplica a la subdivisión heredada; su transporte a una actualización
TPK que cambie esas direcciones se decide componiendo el lector de dicha
actualización.

Existe una formulación hamiltoniana regular de la coordenada dual. En el
plano cotangente con coordenadas \((q,p)\), tome
\(H(q,p)=\psi(p)\). Las ecuaciones son
\(\dot q=\psi'(p)\), \(\dot p=0\). Para velocidades
\(\dot q\in(u_-,u_+)\), la transformación de Legendre da
\[
 L(q,\dot q)=\psi^*(\dot q),\qquad
 \frac{\partial^2L}{\partial\dot q^2}>0.
\]
Este sistema realiza variacionalmente el potencial del flujo de intervalos.
Su ecuación \(\dot p=0\) muestra que el tiempo hamiltoniano es distinto
del parámetro \(s\) que recorre la colección \eqref{eq:convex-partition}.
La evaluación de \(H\) y \(L\) se conserva bajo la subdivisión del
teorema, porque \(\psi\) y \(\psi^*\) permanecen idénticas.

Por su parte, cualquier flujo \(\dot d=F(d)\) admite en una coordenatización local
del simplex una elevación cotangente con
\(\mathscr H(d,p)=p\cdot F(d)\). Esta segunda construcción reproduce
el propio flujo \eqref{eq:replicator} en su base, pero su Hessiano respecto
de \(p\) es nulo. Las dos realizaciones responden a preguntas distintas:
una produce una dualidad convexa regular y la otra transporta una dinámica
de base dada. La distinción preserva el significado de Hamiltoniano y
Lagrangiano al comparar esta geometría con los lectores físicos del corpus.
